% Part: first-order-logic
% Chapter: sequent-calculus
% Section: quantifier-rules

\documentclass[../../../include/open-logic-section]{subfiles}

\begin{document}

\olfileid{fol}{seq}{qrl}

\olsection{Kwantorregels}

\subsection{Regels voor $\lforall$}

\begin{defish}
\Axiom$ !A(t), \Gamma \fCenter \Delta$
\RightLabel{\LeftR{\lforall}}
\UnaryInf$ \lforall[x][!A(x)],\Gamma \fCenter \Delta$
\DisplayProof
\hfill
\Axiom$ \Gamma \fCenter \Delta, !A(a) $
\RightLabel{\RightR{\lforall}}
\UnaryInf$ \Gamma \fCenter \Delta, \lforall[x][!A(x)]$
\DisplayProof
\end{defish}

In \LeftR{\lforall} is $t$ een gesloten term, dat wil zeggen een term
zonder variabelen. In \RightR{\lforall} is $a$~!!a{constant} die
nergens in de onderste sequent van de regel \RightR{\lforall} mag
voorkomen. We noemen $a$ de \emph{eigenvariabele} van de
\RightR{\forall}-afleidingsstap.\footnote{We gebruiken de term
  ``eigenvariabele'' hoewel $a$ in de bovenstaande regel
  !!a{constant} is. Dat heeft historische redenen.}

\subsection{Regels voor $\lexists$}

\begin{defish}
\Axiom$ !A(a), \Gamma \fCenter \Delta $
\RightLabel{\LeftR{\lexists}}
\UnaryInf$ \lexists[x][!A(x)], \Gamma \fCenter \Delta$
\DisplayProof
\hfill
\Axiom$ \Gamma \fCenter \Delta, !A(t) $
\RightLabel{\RightR{\lexists}}
\UnaryInf$ \Gamma \fCenter \Delta, \lexists[x][!A(x)]$
\DisplayProof
\end{defish}

Ook hier is $t$~een gesloten term en is $a$~!!a{constant} die niet in
de onderste sequent van de regel \LeftR{\lexists} voorkomt. We noemen
$a$ de \emph{eigenvariabele} van de \LeftR{\lexists}-afleidingsstap.

De voorwaarde dat een eigenvariabele niet in de onderste sequent van
de \RightR{\lforall}- of \LeftR{\lexists}-afleidingsstap voorkomt,
heet de \emph{eigenvariabelevoorwaarde}.

\begin{explain}
Volgens onze afspraak geven we, als $!A$ !!a{formula} is waarin de
!!{variable}~$x$ vrij voorkomt, dit aan door~$!A(x)$ te schrijven. In dezelfde context is
$!A(t)$ dan een afkorting voor~$\Subst{!A}{t}{x}$. We zouden de regel
\RightR{\lexists} dus ook als volgt kunnen schrijven:
\begin{prooftree}
  \Axiom$\Gamma \fCenter \Delta, \Subst{!A}{t}{x}$
  \RightLabel{\RightR{\lexists}}
  \UnaryInf$\Gamma \fCenter \Delta, \lexists[x][!A]$
\end{prooftree}
Merk op dat $t$ al in~$!A$ mag voorkomen. $!A$~kan bijvoorbeeld
$\Atom{\Obj P}{t,x}$ zijn. De afleiding van $\Gamma \Sequent \Delta,
\lexists[x][\Atom{\Obj P}{t,x}]$ uit~$\Gamma \Sequent \Delta,
\Atom{\Obj P}{t,t}$ is dus een correcte toepassing
van~\RightR{\lexists}: in de premisse mogen een of meer, maar niet
noodzakelijk alle, voorkomens van~$t$ door de gebonden
!!{variable}~$x$ worden ``vervangen''. De eigenvariabelevoorwaarden
voor \RightR{\lforall} en~\LeftR{\lexists} vereisen daarentegen dat de
!!{constant}~$a$ niet in~$!A$ voorkomt. Men kan dus niet op correcte
wijze $\Gamma \Sequent \Delta, \lforall[x][\Atom{\Obj P}{a,x}]$
afleiden uit $\Gamma \Sequent \Delta, \Atom{\Obj P}{a,a}$
met~$\RightR{\lforall}$.
\end{explain}

\begin{explain}
In \RightR{\lexists} en \LeftR{\lforall} gelden voor de term~$t$ geen
verdere beperkingen. Bij de regels \LeftR{\lexists} en
\RightR{\lforall} vereist de eigenvariabelevoorwaarde daarentegen dat
de !!{constant}~$a$ nergens buiten~$!A(a)$ in de bovenste sequent
voorkomt. Deze voorwaarde is nodig om te waarborgen dat het systeem
correct is, dat wil zeggen dat het alleen geldige sequenten
!!{derive}s. Zonder deze voorwaarde zouden de volgende afleidingen zijn
toegestaan:
\begin{prooftree}
  \Axiom$!A(a) \fCenter !A(a)$
  \RightLabel{*\LeftR{\lexists}}
  \UnaryInf$\lexists[x][!A(x)] \fCenter !A(a)$
  \RightLabel{\RightR{\lforall}}
  \UnaryInf$\lexists[x][!A(x)] \fCenter \lforall[x][!A(x)]$
  \DisplayProof\bottomAlignProof
  \qquad
  \Axiom$!A(a) \fCenter !A(a)$
  \RightLabel{*\RightR{\lforall}}
  \UnaryInf$!A(a) \fCenter \lforall[x][!A(x)]$
  \RightLabel{\LeftR{\lexists}}
  \UnaryInf$\lexists[x][!A(x)] \fCenter \lforall[x][!A(x)]$
\end{prooftree}
De sequent $\lexists[x][!A(x)] \Sequent \lforall[x][!A(x)]$ is echter
niet geldig.
\end{explain}
\end{document}
